\documentclass{article}
\usepackage[margin=2cm]{geometry}
\usepackage{amsmath}
\usepackage{graphicx}
\usepackage{booktabs}
\usepackage[utf8]{inputenc}
\usepackage{ragged2e}
\setlength{\emergencystretch}{3em}
\begin{document}
\sloppy

\noindent\hfill\includegraphics[height=1.8em]{assets/figure_header_logo.png}\par\noindent\hrule\vspace{4pt}

with the so-called "boosted structure constants' K given by,.


\begin{equation}
\begin{array}{l l}{{K_{r.a}=g\,\epsilon_{t m n}{\cal L}^{\ell}{}_{r}({\cal L}^{-1}){}_{s}{}^{m}{}_{L}{}^{n}{}_{\alpha}+\lambda\,C_{I J K}{\cal L}^{I}{}_{~r}({\cal L}^{-1}){}_{s}{}^{J}{}_{~}}}\\ {{K_{\alpha\,t t}=g\,\epsilon_{s m n}{\cal L}^{\ell}{}_{\alpha}({\cal L}^{-1})_{I}{}^{m}{}_{t}+\lambda\,C_{M J K}{\cal L}^{M}{}_{~\alpha}({\cal L}^{-1})_{I}{}^{J}{}_{~t}}}\end{array}\qquad\qquad\qquad(1)
\end{equation}


We remind the reader that r, s,t = 1, 2, 3 are obtained from splitting the index  into a 0.
index and an SU(2)R adjoint index. Also appearing in the Lagrangian is Noo, which is the.
00 component of the matrix


\begin{equation}
\mathcal{N}_{\mathrm{AE}}=L_{\mathrm{A}}^{~\alpha}\left(L^{-1}\right)_{\alpha\mathrm{E}}-L_{\mathrm{A}}^{~I}\left(L^{-1}\right)_{L\mathrm{E}}\qquad (2\right)
\end{equation}


\subsection*{0.1 Supersymmetry variations}


theory. In the following section, we will discuss the continuation of this theory to Euclidean
signature, which is complicated by the necessary modification of the symplectic Majorana
condition imposed on the spinor fields. In order to write the fermionic variations, it is first
necessary to introduce a matrix  defined as


\begin{equation}
\gamma^{7}=i\gamma^{0}\gamma^{1}\gamma^{2}\gamma^{3}\gamma^{4}\gamma^{5}\qquad (3)
\end{equation}


and satisfying (y?)2 = -1. With this, the supersymmetry transformations of the fermions
in the Lorentzian case can be given as.


\begin{equation}
\begin{array}{l l}{{\delta\chi_{A}=\frac{i}{2}\gamma^{\mu}\partial_{\mu}\sigma\varepsilon_{A}+N_{A B}\varepsilon^{B}}}\\ {{\delta\lambda_{A}^{I}=i\hat{P}_{\tau^{I}}\sigma_{A B}^{\nu}\partial_{\mu}\phi^{i}\gamma^{\prime}\gamma^{\mu}\varepsilon^{B}-i\hat{P}_{0i}^{I}\epsilon_{A B}\partial_{\mu}\phi^{i}\gamma^{\prime}\gamma^{\mu}\varepsilon^{B}+M_{A B}\varepsilon^{B}}}\\ {{\delta\lambda_{A}^{I}=i\hat{P}_{\tau^{I}\sigma_{A B}^{\nu}\partial_{\mu}\phi^{i}\gamma^{\prime}\gamma^{\mu}\varepsilon^{B}-i\hat{P}_{0i}^{I}\epsilon_{A B}^{I}\epsilon^{\prime}\gamma^{\prime}\epsilon^{B}}}}\end{array}
\end{equation}


where we have defined


\begin{equation}
\begin{array}{l l}{{S_{A B}=\frac{i}{24}[A e^{\sigma}+6m e^{-3\sigma}(L^{-1})_{00}]\varepsilon_{A B}-\frac{i}{8}[B_{t}e^{\sigma}-2m e^{-3\sigma}(L^{-1})_{t0}]\gamma^{7}\sigma_{A B}^{t}}}\\ {{\phantom{\qquad}}}&{{\cdots}}\\ {{N_{A B}=\frac{1}{24}[A e^{\sigma}-18m e^{-3\sigma}(L^{-1})_{t0}]\varepsilon_{A B}+\frac{1}{8}[B_{t}e^{\sigma}+6m e^{-3\sigma}(L^{-1})_{t0}\gamma^{7}\sigma_{A B}^{t}}}\\ {{\phantom{\qquad}}}&{{\cdots}}\end{array}
\end{equation}


In the above, the matrix oAB defined as oAB = orBecA is symmetric in A, B. For more
details, see our previous paper.

\end{document}
